% !TeX program = lualatex
\documentclass[12pt,aspectratio=169]{beamer}
\usepackage[utf8]{inputenc}
\usepackage[T1]{fontenc}
\usepackage{amsmath,amsfonts,amssymb}
\usepackage{graphicx}
\usepackage{xcolor}
\usepackage{hyperref}
\usepackage{anyfontsize}
\usepackage{lmodern}

% ====== EXTREME FONT REDUCTION ======
\renewcommand{\tiny}{\fontsize{4}{5}\selectfont}
\renewcommand{\scriptsize}{\fontsize{5}{6}\selectfont}
\renewcommand{\footnotesize}{\fontsize{6}{7}\selectfont}
\renewcommand{\small}{\fontsize{7}{8}\selectfont}
\renewcommand{\normalsize}{\fontsize{8}{9}\selectfont}
\renewcommand{\large}{\fontsize{9}{10}\selectfont}
\renewcommand{\Large}{\fontsize{10}{11}\selectfont}
\renewcommand{\LARGE}{\fontsize{11}{13}\selectfont}
\renewcommand{\huge}{\fontsize{12}{14}\selectfont}
\renewcommand{\Huge}{\fontsize{14}{17}\selectfont}

% ====== COLOR THEME ======
\definecolor{redcorrect}{RGB}{200,30,30}
\definecolor{bluecorrect}{RGB}{30,80,200}
\definecolor{greencheck}{RGB}{0,150,50}
\definecolor{lightgray}{RGB}{245,245,245}
\definecolor{darkgray}{RGB}{50,50,50}
\definecolor{softyellow}{RGB}{255,248,200}
\definecolor{steppurple}{RGB}{150,50,200}
\definecolor{deepblue}{RGB}{20,60,150}
\definecolor{darkred}{RGB}{150,20,20}

\setbeamercolor{title}{fg=white,bg=redcorrect}
\setbeamercolor{frametitle}{fg=white,bg=redcorrect}
\setbeamercolor{block title}{fg=white,bg=bluecorrect}
\setbeamercolor{block body}{fg=darkgray,bg=lightgray}
\setbeamercolor{normal text}{fg=darkgray}
\usecolortheme{default}

% ====== CUSTOM COMMANDS ======
\newcommand{\correct}[1]{\textcolor{greencheck}{\textbf{\checkmark}} #1}
\newcommand{\keypoint}[1]{\textcolor{deepblue}{\textbf{#1}}}
\newcommand{\hardconcept}[1]{\textcolor{darkred}{\textbf{#1}}}
\newcommand{\tipbox}[1]{\fcolorbox{bluecorrect}{softyellow}{\parbox{0.88\textwidth}{\centering #1}}}
\newcommand{\stepbox}[1]{%
	\begin{center}
		\fcolorbox{darkgray}{white}{%
			\parbox{0.90\textwidth}{\vspace{0.05cm} #1 \vspace{0.05cm}}%
		}
	\end{center}
}

\begin{document}
	
	% ================================================================
	% SLIDE 1: TITLE
	% ================================================================
	\begin{frame}
		\centering
		\vspace{0.8cm}
		{\fontsize{16}{19}\selectfont \textbf{AML / CFT Test Review}} \\[0.15cm]
		{\fontsize{12}{14}\selectfont Advanced Concepts \& Detailed Analysis}\\[0.3cm]
		{\fontsize{9}{11}\selectfont Understanding the \textit{Regulatory Principles} Behind Each Answer}\\[0.3cm]
		{\fontsize{8}{10}\selectfont Compliance Training Department}
		\vspace{0.2cm}
		\rule{4cm}{0.5pt}
	\end{frame}
	
	% ================================================================
	% SLIDE 2: Data Collection — AFC Data Extraction (Q4)
	% ================================================================
	\begin{frame}
		\frametitle{\fontsize{11}{13}\selectfont 1. AFC Data Extraction Challenges}
		\begin{block}{\fontsize{8}{10}\selectfont Subject: Data Collection and Preparation}
			\scriptsize
			Which of the following best describes a challenge during AFC data extraction?
		\end{block}
		\vspace{0.05cm}
		\stepbox{\scriptsize
			\keypoint{Correct Answer:} Data may come in different formats, requiring normalization during integration.
		}
		\vspace{0.05cm}
		\scriptsize
		\hardconcept{Core Concept:} Data extraction involves pulling information from disparate systems — core banking, CRM, trade finance, payments. Each system may use different \textbf{data formats}, \textbf{field names}, and \textbf{data types}. Normalization is the process of transforming raw data into a consistent, usable format. This is a key principle under \textbf{BCBS 239} (Principles for Effective Risk Data Aggregation) and \textbf{FATF Recommendation 15} (New Technologies). Without normalization, monitoring systems cannot properly analyze transactions or identify patterns.
		\vspace{0.05cm}
		\tipbox{\scriptsize \textbf{Key Insight:} Data normalization is not just a technical issue — it has regulatory implications. Inconsistent data formats can lead to \textbf{false negatives} and \textbf{inaccurate reporting}.}
	\end{frame}
	
	% ================================================================
	% SLIDE 3: Data Collection — Data Mining Requirements (Q5)
	% ================================================================
	\begin{frame}
		\frametitle{\fontsize{11}{13}\selectfont 2. Data Mining \& Matching Requirements}
		\begin{block}{\fontsize{8}{10}\selectfont Subject: Data Collection and Preparation}
			\scriptsize
			Which of the following are requirements for effective data mining and matching in AFC controls? (Select Two.)
		\end{block}
		\vspace{0.05cm}
		\stepbox{\scriptsize
			\keypoint{Correct Answers:}
			\par Comparing extracted data against structured and unstructured risk sources.
			\par Applying intelligent filters to reduce false positives during entity resolution.
		}
		\vspace{0.05cm}
		\scriptsize
		\hardconcept{Core Concept:} Data mining in AFC involves \textbf{entity resolution} — connecting data points across multiple sources to identify the same entity (customer, counterparty, or beneficial owner). \textbf{Structured data} includes databases, spreadsheets, and formatted fields. \textbf{Unstructured data} includes emails, documents, and free-text fields. Intelligent filters reduce \textbf{false positives} by applying \textbf{matching thresholds}, \textbf{fuzzy logic}, and \textbf{contextual rules}. This aligns with \textbf{FATF Recommendation 10} (CDD) and \textbf{EU AMLD5} requirements for ongoing monitoring.
		\vspace{0.05cm}
		\tipbox{\scriptsize \textbf{Best Practice:} Implement an \textbf{entity resolution engine} that uses \textbf{machine learning} to improve matching accuracy over time.}
	\end{frame}
	
	% ================================================================
	% SLIDE 4: Data Collection — Effective Integration (Q2)
	% ================================================================
	\begin{frame}
		\frametitle{\fontsize{11}{13}\selectfont 3. Multiple Data Source Integration}
		\begin{block}{\fontsize{8}{10}\selectfont Subject: Data Collection and Preparation}
			\scriptsize
			Which of the following help ensure effective integration of multiple data sources into compliance platforms? (Select Three.)
		\end{block}
		\vspace{0.05cm}
		\stepbox{\scriptsize
			\keypoint{Correct Answers:}
			\par Maintaining a common data dictionary across systems.
			\par Mapping all incoming data to standardized field names.
			\par Aligning control data requirements to system inputs.
		}
		\vspace{0.05cm}
		\scriptsize
		\hardconcept{Core Concept:} Integration success depends on \textbf{data governance}. A \textbf{common data dictionary} ensures all systems use the same definitions (e.g., "customer" is defined consistently). \textbf{Field mapping} ensures that data from source systems is translated correctly into target systems. \textbf{Alignment of control requirements} ensures that the monitoring system receives all necessary data to detect suspicious activity. This is foundational to \textbf{FATF Recommendation 18} (Internal Controls) and \textbf{BCBS 239}.
		\vspace{0.05cm}
		\tipbox{\scriptsize \textbf{Key Principle:} \textbf{Data Governance} is the framework that ensures data is accurate, consistent, and usable across all systems.}
	\end{frame}
	
	% ================================================================
	% SLIDE 5: Data Collection — Internal vs External Sources (Q3)
	% ================================================================
	\begin{frame}
		\frametitle{\fontsize{11}{13}\selectfont 4. Internal vs External Data Sources}
		\begin{block}{\fontsize{8}{10}\selectfont Subject: Data Collection and Preparation}
			\scriptsize
			Which characteristics describe internal versus external sources used during AFC data extraction? (Select Two.)
		\end{block}
		\vspace{0.05cm}
		\stepbox{\scriptsize
			\keypoint{Correct Answers:}
			\par Internal sources reflect observed customer behavior or system-held values.
			\par External sources include sanction or adverse media lists from agencies.
		}
		\vspace{0.05cm}
		\scriptsize
		\hardconcept{Core Concept:} \textbf{Internal data} is generated by the institution's own systems — transaction histories, account balances, customer profiles, onboarding records. \textbf{External data} comes from third-party sources — sanctions lists (OFAC, UN), PEP lists, adverse media databases, and regulatory watchlists. Both are essential for effective AFC monitoring. Internal data provides \textbf{behavioral context}, while external data provides \textbf{risk context}. This dual-source approach is recommended by \textbf{FATF Recommendation 10} and \textbf{EU AMLD5}.
		\vspace{0.05cm}
		\tipbox{\scriptsize \textbf{Real-World Insight:} Many firms use \textbf{data fusion} — combining internal transaction data with external risk intelligence to create a comprehensive risk view.}
	\end{frame}
	
	% ================================================================
	% SLIDE 6: Data Collection — Matching Logic (Q1)
	% ================================================================
	\begin{frame}
		\frametitle{\fontsize{11}{13}\selectfont 5. Effective Data Matching Logic}
		\begin{block}{\fontsize{8}{10}\selectfont Subject: Data Collection and Preparation}
			\scriptsize
			Effective AFC data matching relies heavily on the use of matching logic that can:
		\end{block}
		\vspace{0.05cm}
		\stepbox{\scriptsize
			\keypoint{Correct Answer:} account for slight inconsistencies in names or identifiers.
		}
		\vspace{0.05cm}
		\scriptsize
		\hardconcept{Core Concept:} This is about \textbf{fuzzy matching} and \textbf{phonetic matching} algorithms (e.g., Soundex, Levenshtein distance). Names can appear in different formats (e.g., "Mohammed" vs "Mohamed"), contain typos, or include middle initials. Identifiers (e.g., passport numbers) may have transposed digits. Effective matching logic must account for these variations to avoid missing matches. This is critical for \textbf{sanctions screening} and \textbf{PEP identification}, which are core requirements under \textbf{FATF Recommendation 6} and \textbf{7}.
		\vspace{0.05cm}
		\tipbox{\scriptsize \textbf{Key Insight:} Advanced matching systems use \textbf{machine learning} to continuously improve matching accuracy based on false positive and false negative feedback.}
	\end{frame}
	
	% ================================================================
	% SLIDE 7: Global AFC Standards — Recommendation 18 (Q3)
	% ================================================================
	\begin{frame}
		\frametitle{\fontsize{11}{13}\selectfont 6. Recommendation 18 Implementation}
		\begin{block}{\fontsize{8}{10}\selectfont Subject: Global AFC Standards and Guidance}
			\scriptsize
			A multinational institution is criticized during an audit for deficiencies in Recommendation 18 implementation. What should the auditors have likely focused on?
		\end{block}
		\vspace{0.05cm}
		\stepbox{\scriptsize
			\keypoint{Correct Answer:} Whether group-wide AML programs apply to all branches and subsidiaries.
		}
		\vspace{0.05cm}
		\scriptsize
		\hardconcept{Core Concept:} \textbf{FATF Recommendation 18} requires financial institutions to implement \textbf{group-wide AML/CFT programs}. This includes policies and procedures that apply to all branches and majority-owned subsidiaries. The recommendation also requires that information be shared within the group for AML purposes (subject to data protection laws). The auditors would examine whether the group-wide program is \textbf{consistent} across jurisdictions and whether there is effective \textbf{oversight} of foreign branches.
		\vspace{0.05cm}
		\tipbox{\scriptsize \textbf{Real-World:} This is particularly challenging for multinational banks with operations in countries with weaker AML regimes. The group program must set a \textbf{minimum standard} that applies globally.}
	\end{frame}
	
	% ================================================================
	% SLIDE 8: Global AFC Standards — Three Lines of Defense (Q4)
	% ================================================================
	\begin{frame}
		\frametitle{\fontsize{11}{13}\selectfont 7. Three Lines of Defense}
		\begin{block}{\fontsize{8}{10}\selectfont Subject: Global AFC Standards and Guidance}
			\scriptsize
			A bank is restructuring its AML/CFT framework according to Basel Committee guidelines. Which option correctly describes the three lines of defense in a bank's AML efforts as outlined by the Basel Committee?
		\end{block}
		\vspace{0.05cm}
		\stepbox{\scriptsize
			\keypoint{Correct Answer:} First line: Business units that identify, assess, and control risks; Second line: AML compliance and internal controls; Third line: Internal audit function.
		}
		\vspace{0.05cm}
		\scriptsize
		\hardconcept{Core Concept:} The \textbf{Three Lines of Defense} model is the foundation of AML governance. \textbf{First line}: Operational management owns and manages risk. \textbf{Second line}: Compliance provides oversight and challenge. \textbf{Third line}: Internal audit provides independent assurance. This model is endorsed by the \textbf{Basel Committee on Banking Supervision}, \textbf{FATF}, and \textbf{EU Banking Union}. It ensures \textbf{separation of duties} and \textbf{independent oversight} — critical for effective AML control.
		\vspace{0.05cm}
		\tipbox{\scriptsize \textbf{Key Principle:} The three lines must operate independently but collaboratively. The second line should not replace first-line accountability, and the third line should have direct access to the board.}
	\end{frame}
	
	% ================================================================
	% SLIDE 9: Global AFC Standards — FATF Effectiveness (Q4)
	% ================================================================
	\begin{frame}
		\frametitle{\fontsize{11}{13}\selectfont 8. FATF Effectiveness Assessment}
		\begin{block}{\fontsize{8}{10}\selectfont Subject: Global AFC Standards and Guidance}
			\scriptsize
			Which of the following would FATF assess under the 'effectiveness' component of a mutual evaluation? (Select Three.)
		\end{block}
		\vspace{0.05cm}
		\stepbox{\scriptsize
			\keypoint{Correct Answers:}
			\par Law enforcement use of STRs in convictions.
			\par Actual implementation of AML policies across institutions.
			\par Use of supervisory powers to apply targeted remediation.
		}
		\vspace{0.05cm}
		\scriptsize
		\hardconcept{Core Concept:} FATF mutual evaluations assess both \textbf{technical compliance} (does the law exist?) and \textbf{effectiveness} (is the law producing results?). Effectiveness is measured against \textbf{11 Immediate Outcomes}. The selected options demonstrate:
		\par \textbf{Outcome 7}: STRs used to initiate investigations and lead to convictions.
		\par \textbf{Outcome 4}: AML policies are effectively implemented by institutions.
		\par \textbf{Outcome 3}: Supervisors take action to address non-compliance.
		\vspace{0.05cm}
		\tipbox{\scriptsize \textbf{Key Insight:} Effectiveness is about \textbf{impact}, not just \textbf{process}. Regulators want to see that AML systems actually work.}
	\end{frame}
	
	% ================================================================
	% SLIDE 10: Global AFC Standards — ICAR Support (Q5)
	% ================================================================
	\begin{frame}
		\frametitle{\fontsize{11}{13}\selectfont 9. Basel Institute ICAR Support}
		\begin{block}{\fontsize{8}{10}\selectfont Subject: Global AFC Standards and Guidance}
			\scriptsize
			A regional financial crime task force needs practical assistance with cross-border money laundering cases and wants to enhance their investigative capabilities. Which of the following best describes how the Basel Institute on Governance's International Centre for Asset Recovery (ICAR) could support their needs?
		\end{block}
		\vspace{0.05cm}
		\stepbox{\scriptsize
			\keypoint{Correct Answer:} Providing case mentoring and facilitating international cooperation between jurisdictions.
		}
		\vspace{0.05cm}
		\scriptsize
		\hardconcept{Core Concept:} The \textbf{Basel Institute on Governance}'s ICAR provides \textbf{practical assistance} to law enforcement and prosecution agencies. This includes:
		\par \textbf{Case mentoring}: Hands-on support for complex cases.
		\par \textbf{International cooperation}: Facilitating the exchange of information and evidence between jurisdictions.
		\par \textbf{Asset recovery training}: Technical assistance in tracing and recovering proceeds of crime.
		This aligns with \textbf{FATF Recommendation 38} (Mutual Legal Assistance) and \textbf{UNODC} asset recovery guidelines.
		\vspace{0.05cm}
		\tipbox{\scriptsize \textbf{Real-World:} ICAR has supported cases in the Balkans, Latin America, and Asia-Pacific, helping countries recover billions in stolen assets.}
	\end{frame}
	
	% ================================================================
	% SLIDE 11: Global AFC Standards — Recommendation 36 (Q1)
	% ================================================================
	\begin{frame}
		\frametitle{\fontsize{11}{13}\selectfont 10. Recommendation 36 — Mutual Legal Assistance}
		\begin{block}{\fontsize{8}{10}\selectfont Subject: Global AFC Standards and Guidance}
			\scriptsize
			Which of the following accurate statements regarding FATF Recommendation 36 and its focus on international cooperation? (Select Two.)
		\end{block}
		\vspace{0.05cm}
		\stepbox{\scriptsize
			\keypoint{Correct Answers:}
			\par Countries should provide the widest possible range of mutual legal assistance (MLA).
			\par Authorities may take spontaneous measures to share information.
		}
		\vspace{0.05cm}
		\scriptsize
		\hardconcept{Core Concept:} \textbf{FATF Recommendation 36} requires countries to provide \textbf{mutual legal assistance} in money laundering and terrorist financing investigations. The "widest possible range" means countries should assist in both criminal and civil proceedings. \textbf{Spontaneous information sharing} allows authorities to proactively share information without a formal request, which is critical for fast-moving financial crime cases. Both provisions are essential for \textbf{international cooperation}, which is a key pillar of the FATF framework.
		\vspace{0.05cm}
		\tipbox{\scriptsize \textbf{Key Principle:} MLA is the \textbf{backbone} of cross-border financial crime investigations. Without it, criminals can exploit jurisdictional boundaries.}
	\end{frame}
	
	% ================================================================
	% SLIDE 12: Global AFC Standards — Financial Secrecy Index (Q2)
	% ================================================================
	\begin{frame}
		\frametitle{\fontsize{11}{13}\selectfont 11. Financial Secrecy Index}
		\begin{block}{\fontsize{8}{10}\selectfont Subject: Global AFC Standards and Guidance}
			\scriptsize
			Which of the following might contribute to a jurisdiction's elevated rating according to the Financial Secrecy Index? (Select Three.)
		\end{block}
		\vspace{0.05cm}
		\stepbox{\scriptsize
			\keypoint{Correct Answers:}
			\par Extensive cross-border services to non-residents.
			\par Minimal reporting of beneficial ownership.
			\par High ratio of secrecy score to global scale weight.
		}
		\vspace{0.05cm}
		\scriptsize
		\hardconcept{Core Concept:} The \textbf{Financial Secrecy Index (FSI)} is published by the \textbf{Tax Justice Network}. It measures a jurisdiction's contribution to global financial secrecy. Key contributors include:
		\par \textbf{Cross-border services}: Providing banking and financial services to non-residents enables secrecy.
		\par \textbf{Beneficial ownership reporting}: Minimal reporting allows anonymous ownership.
		\par \textbf{Secrecy score to scale weight ratio}: Indicates that secrecy is disproportionately high relative to the jurisdiction's global financial activity.
		The FSI is used by banks in \textbf{country risk assessments} under \textbf{FATF Recommendation 1}.
		\vspace{0.05cm}
		\tipbox{\scriptsize \textbf{Key Insight:} High FSI ratings are correlated with higher money laundering risk. Banks should apply \textbf{enhanced due diligence} for clients with connections to high-FSI jurisdictions.}
	\end{frame}
	
	% ================================================================
	% SLIDE 13: Use of Guidance — MLAT Challenges (Q1)
	% ================================================================
	\begin{frame}
		\frametitle{\fontsize{11}{13}\selectfont 12. MLAT Implementation Challenges}
		\begin{block}{\fontsize{8}{10}\selectfont Subject: Use of Guidance and AFC Cooperation}
			\scriptsize
			Two countries, A and B, are attempting to establish a mutual legal assistance treaty. Country A has data protection rules preventing sharing of client identities, while Country B bases most financial investigations on beneficial ownership disclosures. What challenge is most likely to arise in implementing the MLAT?
		\end{block}
		\vspace{0.05cm}
		\stepbox{\scriptsize
			\keypoint{Correct Answer:} Evidence exchanged under the MLAT will not be admissible in Country B's legal system.
		}
		\vspace{0.05cm}
		\scriptsize
		\hardconcept{Core Concept:} MLATs are governed by \textbf{mutual legal assistance principles}. A key challenge is \textbf{admissibility of evidence} — evidence gathered in one country must be legally usable in the requesting country's courts. Country A's data protection rules may prevent sharing certain client data (GDPR, PII laws). If Country B's legal system has strict rules on evidence handling, the evidence may be deemed \textbf{inadmissible}. This relates to \textbf{FATF Recommendation 36} and \textbf{UN Convention against Transnational Organized Crime} (UNTOC).
		\vspace{0.05cm}
		\tipbox{\scriptsize \textbf{Real-World:} Many MLATs include provisions for data protection safeguards and chain-of-custody requirements to ensure evidence admissibility.}
	\end{frame}
	
	% ================================================================
	% SLIDE 14: Use of Guidance — PPP Framework (Q2)
	% ================================================================
	\begin{frame}
		\frametitle{\fontsize{11}{13}\selectfont 13. Public-Private Partnerships (PPP) Framework}
		\begin{block}{\fontsize{8}{10}\selectfont Subject: Use of Guidance and AFC Cooperation}
			\scriptsize
			Which of the following best describes the principle behind public-private partnerships (PPPs) in AML efforts?
		\end{block}
		\vspace{0.05cm}
		\stepbox{\scriptsize
			\keypoint{Correct Answer:} Collaboration allows timely information sharing and joint targeting of high-risk actors using shared intelligence.
		}
		\vspace{0.05cm}
		\scriptsize
		\hardconcept{Core Concept:} PPPs are formal arrangements between \textbf{public authorities} (e.g., FIUs, law enforcement) and \textbf{private sector entities} (banks, fintechs). The goal is \textbf{real-time intelligence sharing} to detect and disrupt financial crime. PPPs enable:
		\par \textbf{Timely information sharing}: Faster than traditional channels.
		\par \textbf{Joint targeting}: Combining public intelligence with private sector data.
		\par \textbf{Shared intelligence}: Creating a common understanding of threats.
		This aligns with \textbf{FATF Recommendation 2} (International Cooperation) and \textbf{FATF Recommendation 15} (New Technologies).
		\vspace{0.05cm}
		\tipbox{\scriptsize \textbf{Key Insight:} PPPs are increasingly common in the UK (JMLIT), Australia (Fintel), and Singapore (ACIP). They have proven effective in disrupting money laundering networks.}
	\end{frame}
	
	% ================================================================
	% SLIDE 15: Use of Guidance — PPP Real-Time Intelligence (Q1, second set)
	% ================================================================
	\begin{frame}
		\frametitle{\fontsize{11}{13}\selectfont 14. PPP Real-Time Intelligence Sharing}
		\begin{block}{\fontsize{8}{10}\selectfont Subject: Use of Guidance and AFC Cooperation}
			\scriptsize
			A national financial intelligence unit (FIU) wants to collaborate with financial institutions to disrupt a growing terrorist financing network. It intends to use the public-private partnership (PPP) model for real-time intelligence sharing. Which of the following approaches would be the most aligned with the goals of a PPP framework?
		\end{block}
		\vspace{0.05cm}
		\stepbox{\scriptsize
			\keypoint{Correct Answer:} Establishing a formal platform where regulators and institutions engage in direct two-way communication on specific threats.
		}
		\vspace{0.05cm}
		\scriptsize
		\hardconcept{Core Concept:} Effective PPPs require \textbf{direct, real-time communication}. The platform should facilitate \textbf{two-way information flow} — the FIU shares threat intelligence and institutions share transaction data. Anonymized, after-the-fact reporting is not sufficient for \textbf{disruption} of active networks. This is consistent with the \textbf{FATF Guidance on Public-Private Partnerships} and successful models like the \textbf{UK Joint Money Laundering Intelligence Taskforce (JMLIT)}.
		\vspace{0.05cm}
		\tipbox{\scriptsize \textbf{Real-World:} The JMLIT has been highly effective in disrupting money laundering networks by enabling real-time intelligence sharing between banks and law enforcement.}
	\end{frame}
	
	% ================================================================
	% SLIDE 16: Use of Guidance — Cross-Border Supervision (Q2, second set)
	% ================================================================
	\begin{frame}
		\frametitle{\fontsize{11}{13}\selectfont 15. Cross-Border Regulatory Coordination}
		\begin{block}{\fontsize{8}{10}\selectfont Subject: Use of Guidance and AFC Cooperation}
			\scriptsize
			A cross-border financial institution is supervised by three different regulators across two regions. To maintain comprehensive compliance while reducing administrative burden, what integrated strategy should the regulators apply?
		\end{block}
		\vspace{0.05cm}
		\stepbox{\scriptsize
			\keypoint{Correct Answer:} Regulators should coordinate to clarify scopes, align risk approaches, and execute joint reviews.
		}
		\vspace{0.05cm}
		\scriptsize
		\hardconcept{Core Concept:} Cross-border supervision is governed by the \textbf{Basel Committee's Consolidated Supervision Framework}. The goal is to avoid \textbf{regulatory arbitrage} and \textbf{duplicative examinations}. The integrated strategy includes:
		\par \textbf{Clarifying scopes}: Defining each regulator's jurisdiction.
		\par \textbf{Aligning risk approaches}: Using a common risk methodology.
		\par \textbf{Joint reviews}: Conducting coordinated examinations to reduce burden.
		This aligns with \textbf{FATF Recommendation 2} (International Cooperation) and \textbf{BCBS Principles for the Supervision of Financial Institutions}.
		\vspace{0.05cm}
		\tipbox{\scriptsize \textbf{Key Principle:} \textbf{Home-host supervision} ensures that the lead regulator (home) and host regulators coordinate effectively to maintain comprehensive oversight.}
	\end{frame}
	
	% ================================================================
	% SLIDE 17: Technology — Risk-Based Authentication (Q2)
	% ================================================================
	\begin{frame}
		\frametitle{\fontsize{11}{13}\selectfont 16. Risk-Based Authentication Systems}
		\begin{block}{\fontsize{8}{10}\selectfont Subject: Technology for Customer Onboarding}
			\scriptsize
			How do risk-based authentication systems optimize user experience?
		\end{block}
		\vspace{0.05cm}
		\stepbox{\scriptsize
			\keypoint{Correct Answer:} They prompt additional checks only when user behavior deviates from expectations.
		}
		\vspace{0.05cm}
		\scriptsize
		\hardconcept{Core Concept:} \textbf{Risk-based authentication (RBA)} is a \textbf{dynamic authentication} method that adjusts the level of authentication based on the risk of the transaction. It uses \textbf{behavioral analytics} to establish a baseline of "normal" user behavior. When a user deviates from their normal pattern (e.g., logging in from a new device or location), RBA prompts additional authentication checks. This balances \textbf{security} and \textbf{user experience}, aligning with \textbf{PSD2} (Strong Customer Authentication) and \textbf{FATF Recommendation 15}.
		\vspace{0.05cm}
		\tipbox{\scriptsize \textbf{Key Insight:} RBA reduces friction for legitimate users while adding protective layers for high-risk activities.}
	\end{frame}
	
	% ================================================================
	% SLIDE 18: Technology — Identity Verification Measures (Q3)
	% ================================================================
	\begin{frame}
		\frametitle{\fontsize{11}{13}\selectfont 17. Additional Identity Verification Measures}
		\begin{block}{\fontsize{8}{10}\selectfont Subject: Technology for Customer Onboarding}
			\scriptsize
			A regulator is assessing a fintech startup's customer identity verification program. The firm currently relies solely on device recognition technology for customer authentication. What additional measures could the fintech implement at initial customer onboarding to improve its financial crime controls? (Select Two.)
		\end{block}
		\vspace{0.05cm}
		\stepbox{\scriptsize
			\keypoint{Correct Answers:}
			\par Implementing document verification with biometric authentication.
			\par Integrating with national electronic identity verification systems.
		}
		\vspace{0.05cm}
		\scriptsize
		\hardconcept{Core Concept:} Device recognition alone is insufficient for \textbf{Customer Due Diligence}. Additional measures include:
		\par \textbf{Document verification}: Scanning and validating government-issued IDs (passport, driver's license).
		\par \textbf{Biometric authentication}: Using facial recognition or fingerprint to verify the user is the person on the ID.
		\par \textbf{National e-ID integration}: Connecting to government identity systems (e.g., EU eIDAS, India Aadhaar).
		These measures align with \textbf{FATF Recommendation 10} (CDD) and \textbf{EU AMLD5} requirements for enhanced identity verification.
		\vspace{0.05cm}
		\tipbox{\scriptsize \textbf{Key Principle:} \textbf{Multi-factor authentication} (MFA) is essential for robust identity verification. Device recognition should be just one layer.}
	\end{frame}
	
	% ================================================================
	% SLIDE 19: Technology — AML System Tuning (Q5)
	% ================================================================
	\begin{frame}
		\frametitle{\fontsize{11}{13}\selectfont 18. AML System Post-Implementation Tuning}
		\begin{block}{\fontsize{8}{10}\selectfont Subject: Technology for AFC Compliance}
			\scriptsize
			During post-implementation review, a newly integrated AML system shows a significant increase in alert volume, many of which are false positives. What is the most appropriate next step?
		\end{block}
		\vspace{0.05cm}
		\stepbox{\scriptsize
			\keypoint{Correct Answer:} Adjust alert parameters after conducting a tuning exercise informed by alert productivity metrics and typology validation.
		}
		\vspace{0.05cm}
		\scriptsize
		\hardconcept{Core Concept:} System tuning is an iterative process that must be done \textbf{systematically}. The approach should be:
		\par \textbf{Alert productivity metrics}: Analyze what percentage of alerts result in SAR filings.
		\par \textbf{Typology validation}: Test scenarios against known money laundering methods.
		\par \textbf{Parameter adjustment}: Refine thresholds based on analysis.
		This aligns with \textbf{FATF Recommendation 15} (New Technologies) and \textbf{SR 11-7} (Model Risk Management). Tuning is essential to reduce \textbf{false positives} and maintain operational efficiency.
		\vspace{0.05cm}
		\tipbox{\scriptsize \textbf{Best Practice:} Conduct tuning exercises regularly (at least annually) and after any major change to the system or business environment.}
	\end{frame}
	
	% ================================================================
	% SLIDE 20: Technology — Risk-Based Resource Prioritization (Q6)
	% ================================================================
	\begin{frame}
		\frametitle{\fontsize{11}{13}\selectfont 19. Risk-Based Resource Prioritization}
		\begin{block}{\fontsize{8}{10}\selectfont Subject: Technology for AFC Compliance}
			\scriptsize
			Which of the following actions are consistent with a risk-based approach to prioritizing AFC resources? (Select Three.)
		\end{block}
		\vspace{0.05cm}
		\stepbox{\scriptsize
			\keypoint{Correct Answers:}
			\par Regularly assessing emerging regional threats to adjust mitigation focus.
			\par Prioritizing compliance work that directly mitigates high-risk typologies.
			\par Maintaining flexibility in staffing to respond to urgent tasks.
		}
		\vspace{0.05cm}
		\scriptsize
		\hardconcept{Core Concept:} The \textbf{risk-based approach} is the foundation of modern AML frameworks. It requires:
		\par \textbf{Regular threat assessment}: Monitoring emerging risks (e.g., new money laundering methods).
		\par \textbf{Direct mitigation focus}: Allocating resources to the highest-risk areas.
		\par \textbf{Staffing flexibility}: Having the ability to redeploy staff to urgent issues.
		This aligns with \textbf{FATF Recommendation 1} and the \textbf{Basel Committee's Risk-Based Supervision} principles.
		\vspace{0.05cm}
		\tipbox{\scriptsize \textbf{Key Principle:} The risk-based approach is \textbf{dynamic}, not static. Resources should be allocated based on the current risk landscape.}
	\end{frame}
	
	% ================================================================
	% SLIDE 21: Technology — Behavioral Pattern Integration (Q3)
	% ================================================================
	\begin{frame}
		\frametitle{\fontsize{11}{13}\selectfont 20. Integrating Digital Wallet Patterns}
		\begin{block}{\fontsize{8}{10}\selectfont Subject: Technology for Ongoing Monitoring and Investigations}
			\scriptsize
			A large transaction monitoring system generates alerts for high-value transfers, structured transactions, and velocity anomalies. Compliance officers now want to include behavior patterns involving digital wallets. What is the primary benefit of integrating this pattern?
		\end{block}
		\vspace{0.05cm}
		\stepbox{\scriptsize
			\keypoint{Correct Answer:} It accounts for emerging laundering techniques.
		}
		\vspace{0.05cm}
		\scriptsize
		\hardconcept{Core Concept:} Money laundering methods evolve rapidly. \textbf{Digital wallets} are increasingly used in \textbf{cryptocurrency laundering} and \textbf{structuring}. Integrating digital wallet patterns ensures the monitoring system detects \textbf{emerging typologies} — new ways criminals launder money. This is aligned with \textbf{FATF Recommendation 15} (New Technologies) and the \textbf{FATF Report on Virtual Assets} (VA and VASP recommendations).
		\vspace{0.05cm}
		\tipbox{\scriptsize \textbf{Key Insight:} Monitoring systems must be \textbf{adaptive} to keep pace with criminal innovation. Regular threat assessments should inform system updates.}
	\end{frame}
	
	% ================================================================
	% SLIDE 22: Technology — Batch Screening (Q4)
	% ================================================================
	\begin{frame}
		\frametitle{\fontsize{11}{13}\selectfont 21. Batch Screening}
		\begin{block}{\fontsize{8}{10}\selectfont Subject: Technology for Ongoing Monitoring and Investigations}
			\scriptsize
			Which of the following are true about batch screening? (Select Two.)
		\end{block}
		\vspace{0.05cm}
		\stepbox{\scriptsize
			\keypoint{Correct Answers:}
			\par It helps reveal changing risks within the existing customer base.
			\par It is designed to handle large volumes of data.
		}
		\vspace{0.05cm}
		\scriptsize
		\hardconcept{Core Concept:} \textbf{Batch screening} processes data in bulk — typically the entire customer database — against sanctions lists, PEP lists, and other watchlists. It is:
		\par \textbf{Volume-oriented}: Designed to handle millions of records.
		\par \textbf{Risk-aware}: Detects existing customers who have become risky (e.g., newly designated sanctions).
		\par \textbf{Periodic}: Runs on a schedule (daily, weekly, monthly).
		This is distinct from \textbf{real-time screening} and is a requirement under \textbf{FATF Recommendation 10} (CDD) for ongoing monitoring of existing customers.
		\vspace{0.05cm}
		\tipbox{\scriptsize \textbf{Key Insight:} Batch screening is essential for detecting \textbf{changes in customer risk status} — such as new PEP designations or sanctions updates.}
	\end{frame}
	
	% ================================================================
	% SLIDE 23: Technology — Risk Scoring Engine Results (Q4, second set)
	% ================================================================
	\begin{frame}
		\frametitle{\fontsize{11}{13}\selectfont 22. Customer Risk Scoring Engine}
		\begin{block}{\fontsize{8}{10}\selectfont Subject: Technology for Ongoing Monitoring and Investigations}
			\scriptsize
			An organization is piloting a customer-risk scoring engine that combines internal and external data. After the first quarter, high-risk clients that were previously undetected are flagged. How should this result be interpreted?
		\end{block}
		\vspace{0.05cm}
		\stepbox{\scriptsize
			\keypoint{Correct Answer:} This may indicate successful performance by identifying missed risk exposure.
		}
		\vspace{0.05cm}
		\scriptsize
		\hardconcept{Core Concept:} A \textbf{risk scoring engine} uses \textbf{predictive analytics} to assign risk scores to customers. The ability to flag previously undetected high-risk clients suggests that:
		\par The engine is \textbf{identifying risk} that was missed by previous methods.
		\par The \textbf{data integration} is effective — combining internal and external data reveals a fuller picture.
		\par The \textbf{algorithm} is working as intended.
		This aligns with \textbf{FATF Recommendation 1} (risk-based approach) and \textbf{BCBS 239} (data aggregation).
		\vspace{0.05cm}
		\tipbox{\scriptsize \textbf{Key Insight:} The goal of a risk scoring engine is to \textbf{identify previously unknown or underestimated risk}. Success is measured by its ability to detect overlooked risk.}
	\end{frame}
	
	% ================================================================
	% SLIDE 24: Technology — Real-Time vs Batch Screening (Q5, second set)
	% ================================================================
	\begin{frame}
		\frametitle{\fontsize{11}{13}\selectfont 23. Real-Time vs Batch Screening}
		\begin{block}{\fontsize{8}{10}\selectfont Subject: Technology for Ongoing Monitoring and Investigations}
			\scriptsize
			Which of the following best describes the difference in focus between real-time screening and batch screening against sanctions lists or other watchlists? (Select Two.)
		\end{block}
		\vspace{0.05cm}
		\stepbox{\scriptsize
			\keypoint{Correct Answers:}
			\par Real-time screening evaluates individual customers or transactions at the point of interaction.
			\par Batch screening processes the entire customer database periodically.
			\par Real-time screening prevents high-risk individuals from onboarding.
			\par Batch screening identifies existing customers whose risk status has changed.
		}
		\vspace{0.05cm}
		\scriptsize
		\hardconcept{Core Concept:} \textbf{Real-time screening} is used at the \textbf{point of interaction} — onboarding, transactions, and payments. Its goal is to \textbf{prevent} prohibited activities. \textbf{Batch screening} is used to \textbf{detect} changes in risk status for existing customers. Both are required by \textbf{FATF Recommendation 10} (CDD) and \textbf{Recommendation 7} (Sanctions).
		\vspace{0.05cm}
		\tipbox{\scriptsize \textbf{Key Principle:} Real-time and batch screening are \textbf{complementary}, not substitute. Both are necessary for effective AML controls.}
	\end{frame}
	
	% ================================================================
	% SLIDE 25: Money Laundering — Environmental Crimes (Q4)
	% ================================================================
	\begin{frame}
		\frametitle{\fontsize{11}{13}\selectfont 24. Environmental Crimes \& Money Laundering}
		\begin{block}{\fontsize{8}{10}\selectfont Subject: Money Laundering and Financial Crime}
			\scriptsize
			Why are environmental crimes increasingly associated with money laundering risk?
		\end{block}
		\vspace{0.05cm}
		\stepbox{\scriptsize
			\keypoint{Correct Answer:} They often generate high-volume cash proceeds that require integration.
		}
		\vspace{0.05cm}
		\scriptsize
		\hardconcept{Core Concept:} Environmental crimes — illegal logging, mining, wildlife trafficking, and waste dumping — are \textbf{highly profitable}. They generate \textbf{large amounts of cash} that need to be integrated into the financial system. This makes them a significant source of \textbf{proceeds of crime}. The \textbf{FATF} has issued specific guidance on \textbf{Environmental Crime and Money Laundering}, recognizing the link between environmental crime and financial crime. This is also addressed by \textbf{EU AMLD5} and the \textbf{UN Framework Convention on Climate Change (UNFCCC)}.
		\vspace{0.05cm}
		\tipbox{\scriptsize \textbf{Key Insight:} Financial institutions should consider environmental crime risk as part of their \textbf{Country Risk Assessment} and \textbf{Customer Due Diligence} processes.}
	\end{frame}
	
	% ================================================================
	% SLIDE 26: Money Laundering — Proactive Financial Crime Prevention (Q5)
	% ================================================================
	\begin{frame}
		\frametitle{\fontsize{11}{13}\selectfont 25. Proactive Financial Crime Prevention}
		\begin{block}{\fontsize{8}{10}\selectfont Subject: Money Laundering and Financial Crime}
			\scriptsize
			Which of the following best illustrates an institution's proactive approach to preventing financial crime?
		\end{block}
		\vspace{0.05cm}
		\stepbox{\scriptsize
			\keypoint{Correct Answer:} Utilizing AI for real-time fraud detection across all customer accounts.
		}
		\vspace{0.05cm}
		\scriptsize
		\hardconcept{Core Concept:} \textbf{Proactive} AML means taking \textbf{preventative measures} before financial crime occurs. \textbf{AI-powered fraud detection} is \textbf{proactive} because it:
		\par \textbf{Monitors in real-time}: Detects suspicious activity immediately.
		\par \textbf{Uses predictive analytics}: Identifies patterns that may indicate future risk.
		\par \textbf{Adapts to new methods}: Machine learning models improve over time.
		This is contrasted with \textbf{reactive} approaches (e.g., responding to SARs after the fact). Proactivity is emphasized in \textbf{FATF Recommendation 15} and the \textbf{Basel Committee's Sound Management of Operational Risk}.
		\vspace{0.05cm}
		\tipbox{\scriptsize \textbf{Key Principle:} Proactive AML is about \textbf{stopping crime before it happens}. AI and machine learning are key enablers.}
	\end{frame}
	
	% ================================================================
	% SLIDE 27: Money Laundering — Cyber-Enabled Crime (Q6)
	% ================================================================
	\begin{frame}
		\frametitle{\fontsize{11}{13}\selectfont 26. Cyber-Enabled vs Traditional Fraud}
		\begin{block}{\fontsize{8}{10}\selectfont Subject: Money Laundering and Financial Crime}
			\scriptsize
			Which scenario most suggests cyber-enabled crime rather than traditional fraud?
		\end{block}
		\vspace{0.05cm}
		\stepbox{\scriptsize
			\keypoint{Correct Answer:} A fraudster uses phishing emails to collect customer banking info.
		}
		\vspace{0.05cm}
		\scriptsize
		\hardconcept{Core Concept:} \textbf{Cyber-enabled crime} uses technology to commit traditional crimes. \textbf{Phishing} is a classic example — it uses digital communication to \textbf{deceive} victims into revealing sensitive information. Key characteristics:
		\par \textbf{Digital deception}: Uses email, SMS, or social media.
		\par \textbf{Remote execution}: The fraudster never meets the victim.
		\par \textbf{Large-scale targeting}: Can reach millions of potential victims.
		This is distinct from \textbf{traditional fraud} (e.g., check fraud, physical theft). Cyber-enabled crime is addressed by \textbf{FATF} and \textbf{EU Cybersecurity Act} guidance.
		\vspace{0.05cm}
		\tipbox{\scriptsize \textbf{Key Insight:} Banks must implement \textbf{cybersecurity controls} (email filtering, phishing detection) and train employees to recognize \textbf{social engineering} techniques.}
	\end{frame}
	
	% ================================================================
	% SLIDE 28: SUMMARY
	% ================================================================
	\begin{frame}
		\frametitle{\fontsize{11}{13}\selectfont SUMMARY: Advanced Concepts Matrix}
		\scriptsize
		\begin{block}{\fontsize{9}{11}\selectfont Regulatory Principles \& Their Application:}
			\vspace{0.05cm}
			\scriptsize
			\begin{tabular}{|p{0.35\textwidth}|p{0.60\textwidth}|}
				\hline
				\keypoint{Concept} & \keypoint{Application} \\
				\hline
				FATF Rec 10 (CDD) & EDD for high-risk clients; UBO identification \\
				\hline
				FATF Rec 15 (Tech) & Centralized monitoring; data quality controls \\
				\hline
				FATF Rec 18 (Internal Controls) & Group-wide programs; role-specific training \\
				\hline
				FATF Rec 36 (MLA) & Cross-border evidence sharing; spontaneous disclosure \\
				\hline
				Three Lines of Defense & First line enforcement; second line oversight \\
				\hline
				Model Risk Mgmt (SR 11-7) & Change management for monitoring systems \\
				\hline
				BCBS 239 (Data Aggregation) & Data quality; common data dictionaries \\
				\hline
				Financial Secrecy Index & Country risk assessment for due diligence \\
				\hline
				Public-Private Partnerships & Real-time intelligence sharing; joint targeting \\
				\hline
				Risk-Based Authentication & Dynamic authentication; behavioral analytics \\
				\hline
				Batch vs Real-Time Screening & Detection vs prevention; complementary \\
				\hline
				Environmental Crime & Cash-intensive; integration risk \\
				\hline
				Cyber-Enabled Crime & Phishing; digital deception \\
				\hline
			\end{tabular}
			\vspace{0.05cm}
		\end{block}
		\vspace{0.05cm}
		\tipbox{\scriptsize \textbf{FINAL LESSON:} Every question tests the ability to apply \textbf{regulatory principles} to real-world scenarios. Understand the \textit{why}, and the \textit{what} becomes obvious.}
	\end{frame}
	
	% ================================================================
	% SLIDE 29: THANK YOU
	% ================================================================
	\begin{frame}
		\centering
		\vspace{2.5cm}
		{\fontsize{16}{19}\selectfont \textbf{Thank You}} \\[0.2cm]
		{\fontsize{11}{13}\selectfont Keep learning. Keep growing.} \\[0.15cm]
		\rule{3cm}{0.5pt} \\[0.15cm]
		{\fontsize{8}{10}\selectfont \textit{Compliance is not about perfection — it is about continuous improvement and understanding.}}
		\vspace{0.3cm}
		{\fontsize{7}{9}\selectfont \textcolor{darkgray}{End of Advanced Review}}
	\end{frame}
	
\end{document}